%%=============================================================================
%% Conclusie
%%=============================================================================

\chapter{\IfLanguageName{dutch}{Conclusie}{Conclusion}}%
\label{ch:conclusie}

% TODO: Trek een duidelijke conclusie, in de vorm van een antwoord op de
% onderzoeksvra(a)g(en). Wat was jouw bijdrage aan het onderzoeksdomein en
% hoe biedt dit meerwaarde aan het vakgebied/doelgroep? 
% Reflecteer kritisch over het resultaat. In Engelse teksten wordt deze sectie
% ``Discussion'' genoemd. Had je deze uitkomst verwacht? Zijn er zaken die nog
% niet duidelijk zijn?
% Heeft het onderzoek geleid tot nieuwe vragen die uitnodigen tot verder 
% onderzoek?

Dit onderzoek had als hoofddoel te evalueren of een in-house Master Data Management (MDM)-oplossing, gebouwd binnen het Microsoft Fabric-ecosysteem, een haalbaar en schaalbaar alternatief biedt voor gespecialiseerde commerciële pakketten zoals Profisee. Aan de hand van een Proof of Concept (PoC) is de gefragmenteerde artikel- en leveranciersdata van drie IPCOM-entiteiten (Oostenrijk, Zwitserland en Noorwegen) geëxtraheerd, geharmoniseerd en succesvol ontdubbeld. 

Uit de resultaten blijkt eenduidig dat de in-house ontwikkelde pijplijn uiterst effectief is. De combinatie van een deterministische waterval-logica en probabilistische fuzzy matching (via de Levenshtein-afstand) resulteerde in een deduplicatiegraad van 33,4\%. Hiermee benadert de Fabric-oplossing het niveau van het Profisee-platform (35,1\%) verrassend dicht. De meerwaarde van dit onderzoek voor IPCOM is dan ook aanzienlijk: het toont aan dat datakwaliteit en consolidatie fundamenteel verbeterd kunnen worden zonder de data buiten de eigen vertrouwde cloud-omgeving te verplaatsen en zonder bijkomende licentiekosten voor externe platformen.

\section*{Kritische Reflectie (Discussie)}

Wanneer we kritisch reflecteren op deze uitkomst, overtreft de PoC de initiële verwachtingen op technisch vlak. Vooraf werd aangenomen dat een gespecialiseerde tool als Profisee, met zijn geavanceerde ingebouwde algoritmen, een veel grotere voorsprong zou hebben in het identificeren van complexe, contextuele matches. Dat het verschil slechts marginaal was, bewijst de kracht van PySpark in combinatie met een goed doordacht canonical model en zorgvuldige data cleansing. 

Toch moeten er een aantal belangrijke kanttekeningen worden geplaatst. De grootste beperking van de opgezette Fabric-architectuur is de toegankelijkheid voor niet-technische gebruikers. Waar Profisee beschikt over een intuïtieve, \textit{low-code} interface waarin \textit{data stewards} handmatig randgevallen kunnen beoordelen en corrigeren, vereist de huidige PoC kennis van Python en data engineering. Voor elke wijziging in de matching- of survivorship-regels moet de onderliggende code worden aangepast. Het platform functioneert dus uitstekend als dataverwerkingsmotor, maar mist momenteel een gebruiksvriendelijke beheerlaag.

Daarnaast bleek het bepalen van de drempelwaarde (threshold) voor de fuzzy matching een delicaat proces. Een te lage waarde leidde snel tot \textit{false positives}, terwijl een te hoge waarde de effectiviteit van het algoritme tenietdeed. Het balanceren van deze logica vergt continue monitoring naarmate de dataset groeit.

\section*{Aanbevelingen en Toekomstig Onderzoek}

Dit onderzoek legt een solide technische basis voor MDM binnen Microsoft Fabric, maar opent tegelijkertijd de deur naar nieuwe vraagstukken die verder onderzoek vereisen. 

Ten eerste wordt aanbevolen om te onderzoeken hoe de ontbrekende beheerlaag (de \textit{data steward interface}) binnen het Microsoft-ecosysteem kan worden opgevangen. Een mogelijke piste is de integratie van Power Apps bovenop de Delta Tables in de gold-laag van het Fabric datawarehouse. Hiermee zou een interface op maat gebouwd kunnen worden, zodat domeinexperts datakwaliteitsproblemen kunnen oplossen zonder in de code te hoeven ingrijpen.

Ten tweede nodigt de complexiteit van internationale data uit tot verder onderzoek op het gebied van taalbarrières. Aangezien de entiteiten binnen IPCOM verschillende talen hanteren (zoals Duits, Noors en Nederlands), was het matchen van artikelnamen nu nog deels beperkt door taalkundige verschillen. Het integreren van een vertaalservice (bijvoorbeeld Azure AI Translator) in de harmonisatiefase, om alle omschrijvingen eerst naar het Engels te vertalen alvorens de fuzzy matching toe te passen, zou de deduplicatiegraad potentieel nog verder kunnen verhogen.

Tot slot adviseert dit onderzoek IPCOM om de PoC op termijn stapsgewijs te schalen naar de overige Europese entiteiten. Door de robuustheid van de Lakehouse-architectuur en de rekenkracht van PySpark is de technische fundering daarvoor ruimschoots aanwezig.